\documentclass[12pt,letterpaper]{article}

\usepackage[T1]{fontenc}
\usepackage[utf8]{inputenc}
\usepackage{lmodern}
\usepackage[margin=1in]{geometry}
\usepackage{setspace}
\usepackage{microtype}
\usepackage{amsmath,amssymb,amsthm}
\usepackage{hyperref}
\usepackage[capitalise,nameinlink,noabbrev]{cleveref}
\usepackage{fancyhdr}

\setstretch{1.08}
\setlength{\headheight}{14pt}
\pagestyle{fancy}
\fancyhf{}
\lhead{For the People, By Force}
\rhead{J. Christian McKinley}
\cfoot{\thepage}

\hypersetup{
    colorlinks=true,
    linkcolor=blue,
    citecolor=blue,
    urlcolor=blue
}

\newtheorem{proposition}{Proposition}[section]
\newtheorem{definition}[proposition]{Definition}
\newtheorem{remark}[proposition]{Remark}

\title{For the People, By Force: How Communism and Socialism Use the Same Moral Argument.}
\author{J. Christian McKinley}
\date{22 April 2026}

\begin{document}
\maketitle

\begin{abstract}
Communism achieved through war and socialism achieved through redistribution are commonly
treated as different political types, one revolutionary and one humane. This paper argues
that they use the same justificatory structure. In both cases, an existing order is declared
morally illegitimate, coercion is recoded as justice, seizure is presented as service to the
people, and resistance is treated as vice. The difference is not the argument but the style
of presentation. Communism states the force element more openly and often reaches
concentrated rule more quickly. Socialism advances more gradually under the language of
fairness. Yet both rely on the same moral conversion: compelled surrender is redescribed
as public virtue. The result is movement away from democracy and toward autocracy.
\end{abstract}

\section{Introduction}

Communism says that force is morally necessary to overthrow an unjust order for the people.
Socialism says that compelled redistribution is morally necessary to correct unjust outcomes
for the people. These are usually treated as different arguments because one sounds militant
and the other sounds humane. The difference is largely rhetorical.

The claim of this paper is direct. Communism achieved through war and socialism achieved
through redistribution share the same argumentative skeleton. Both begin by condemning an
existing arrangement as morally defective. Both transform coercion into duty. Both present
seizure as service to the people. Both narrow the legitimacy of refusal. Both expand authority
in the name of justice. The difference is tempo and tone, not structure.

\section{Definitions}

\begin{definition}[Democracy]
Democracy is rule through elected leadership under conditions in which political turnover
remains meaningful, private refusal retains real scope, and the state does not treat ordinary
unequal outcomes as standing warrants for revision by force.
\end{definition}

\begin{definition}[Autocracy]
Autocracy is concentrated rule in which political authority expands beyond the protection of
common order and increasingly claims power to direct, override, or seize private life in the
name of a superior end.
\end{definition}

\begin{definition}[Seizure]
Seizure is compelled surrender of property, production, control, or political authority under
threat of force or sanction, whether presented openly as confiscation or softly as lawful
correction.
\end{definition}

\begin{definition}[Moral laundering]
Moral laundering is the redescription of coercive action as justice, fairness, compassion, or
public duty, such that domination is presented as benevolence.
\end{definition}

\section{The Shared Sequence}

The common structure is simple.

First, an existing order is declared morally illegitimate. In communism, the target is class
rule requiring overthrow. In socialism, the target is unequal outcome requiring correction.

Second, coercion is recoded as justice. In communism, revolutionary force becomes liberation.
In socialism, compelled transfer becomes fairness.

Third, seizure is presented as service to the people. The doctrine does not describe itself as
taking for rulers. It describes itself as acting for a morally privileged collective.

Fourth, resistance is recoded as vice. In communism, opposition becomes sabotage or defense
of oppression. In socialism, opposition becomes selfishness, greed, or indifference to the
vulnerable.

Fifth, authority expands without a clear stopping point. If the first coercive correction is
justified, the next one is easier to justify. If ordinary liberty already failed the moral test,
further intervention appears not as excess but as consistency.

\section{The Common Error of Surface Classification}

A standard political error is to classify doctrines by tone rather than by what they authorize.
Under this error, communism is grouped with overt coercion while socialism is grouped with
humane reform. One sounds harsh. The other sounds compassionate. From that contrast,
many conclude that the two differ in kind.

This is false. A doctrine is not meaningfully classified by whether it says ``war'' or
``fairness.'' It is classified by whether it authorizes compelled surrender under a morally
privileged public end. By that measure, communism and socialism are closer than their
surface language suggests.

Communism is clearer about seizure of command. Socialism is clearer about seizure of output.
Both justify taking what is not voluntarily yielded, and both sanctify that act by appeal to
the people.

\section{Why Resistance Must Be Morally Degraded}

No coercive political project can stabilize while treating refusal as fully legitimate. The
refuser must first be recoded as morally defective.

In communism, resistance is readily treated as counterrevolution, class defense, or sabotage.
In socialism, resistance is treated as selfishness, greed, or refusal to care for others. These
labels differ in rhetoric but not in function. They turn a dissenter into an offender against
justice.

That step matters because it makes force easier to defend. Coercion no longer appears as an
aggression against a rights-bearing citizen. It appears as deserved correction against a person
standing in the way of the people.

\section{Tempo Differs; Structure Does Not}

The obvious objection is that communism is still worse because it uses war, repression, and
overt revolutionary violence, while socialism often works through elections, legislation, and
gradual policy expansion. That objection is true in degree and false in structure.

Yes, the communist route is usually faster, harsher, and more candid about force. Yes, the
socialist route is usually slower, softer in language, and more likely to preserve democratic
forms for longer. But those differences do not alter the underlying pattern. One path says,
``Take power by force for the people.'' The other says, ``Expand compelled surrender for
fairness and the people.'' In both, coercion is justified by moral necessity on behalf of a
collective beneficiary.

The difference is therefore one of style and tempo, not skeleton.

\section{Why the Socialist Variant Is Easier to Miss}

The communist route alarms because its force element is visible. Revolutionary language and
direct command announce themselves. The socialist route is more concealing because it wraps
coercion in humane language. It says not that power must be seized, but that fairness must
be achieved. It says not that opponents must be crushed, but that selfishness must not block
justice.

For that reason, socialism can normalize compelled surrender while continuing to speak as
the critic of domination. Its coercion appears softened by moral language. Its expansion of
authority appears compassionate rather than autocratic. Its path toward concentrated rule is
therefore easier to mistake for democracy improved rather than democracy thinned.

\section{Comparative Proposition}

\begin{proposition}
Communism achieved through war and socialism achieved through redistribution employ the
same basic justificatory structure. Each declares an existing order morally illegitimate, recodes
coercion as justice, presents seizure as service to the people, degrades resistance as vice, and
thereby expands political authority toward autocracy rather than democracy.
\end{proposition}

\begin{proof}
In communism, the existing order is treated as unjust class domination. Forceful overthrow is
therefore justified as moral duty. Seizure of control is framed as liberation for the people.
Resistance is redescribed as sabotage or defense of oppression. Authority centralizes because
the project claims a superior moral mission overriding ordinary limits.

In socialism, unequal outcomes produced by free competition are treated as moral defects
requiring correction. Compelled redistribution is therefore justified as fairness. Seizure of
private output is framed as justice for the people. Resistance is redescribed as selfishness or
lack of compassion. Authority expands because correction of unequal outcomes becomes a
standing moral task.

Thus the rhetorical surface differs, but the operative sequence is the same: moral indictment,
justified coercion, sanctified seizure, degraded refusal, and enlarged rule. Because this
sequence narrows private independence and broadens force-backed political supervision, both
move structurally away from democracy and toward autocracy.
\end{proof}

\section{The Real Contrast}

The real contrast is not between a violent doctrine and a humane doctrine. It is between
openly admitted coercion and morally laundered coercion.

Communism says:
\begin{quote}
Take power by force for the people.
\end{quote}

Socialism says:
\begin{quote}
Take output by law for fairness and the people.
\end{quote}

One sounds revolutionary. The other sounds compassionate. Both involve compelled surrender
justified by collective moral purpose. The difference is not whether force is present. The
difference is whether force is advertised plainly or softened through prestige language.

\section{Conclusion}

Communism achieved through war and socialism achieved through redistribution use the same
moral argument. Both condemn an existing order as illegitimate. Both transform coercion
into supposed justice. Both present seizure as service to the people. Both treat resistance as
moral defect. Both thereby normalize expansion of concentrated political authority.

The difference between them is real but secondary. Communism is more explicit, more
militant, and often faster in its route to concentrated rule. Socialism is softer in language,
more gradual in its advance, and more likely to preserve democratic forms for longer. But
these differences do not change the grammar. In both cases, compelled surrender is sold as
public virtue.

The relevant question, then, is not whether a doctrine speaks in the language of war or
fairness. The relevant question is whether it authorizes seizure under force while claiming to
act for the people. Where that structure appears, democracy is already being thinned. The
road may be loud or quiet. The destination is autocracy.

\section*{Related Reading}

\begin{enumerate}
    \item J. Christian McKinley, \emph{Taxation as Legislated Crime via Majority Hegemony}, Zenodo (2026). DOI: \href{https://doi.org/10.5281/zenodo.18530711}{10.5281/zenodo.18530711}.
    \item J. Christian McKinley, \emph{Destructive Kindness: How Redistribution Inverts Incentives and Shrinks Production}, Zenodo (2025). DOI: \href{https://doi.org/10.5281/zenodo.17773660}{10.5281/zenodo.17773660}.
\end{enumerate}

\end{document}